% !TEX root = ../../main.tex

\section{النمذجة وإدخال البارامترات \el{(Practical Implementation)}}

يركز هذا القسم على الجانب التطبيقي لبناء مرشح التمرير المنخفض \el{(\textbf{Low Pass Filter})} داخل بيئة \el{\textbf{Simulink}}، حيث تم تعديل الموديل السابق ليتناسب مع البنية الجديدة للمرشح التي تتألف من ملفات تسلسلية ومخمدات \el{(RC)} تفرعية.

% --- 3.1 بناء المخطط الصندوقي ---
\subsection{بناء المخطط الصندوقي \el{(Circuit Topology)}}

لتحقيق استجابة "تمرير منخفض"، تم إدراج ملفات خنق \el{(Inductors)} على التوالي مع خطوط التغذية الثلاثة الواصلة إلى المقوم. بعد هذه الملفات مباشرة، تم توصيل فروع مكونة من مقاومة ومكثف على التوالي موصلة نجمياً \el{(Wye-connected)} إلى الأرضي. تضمن هذه البنية إجبار التيارات التوافقية ذات التردد العالي على المرور عبر المسار التفرعي إلى الأرضي، بينما يمر التردد الأساسي \el{($50 Hz$)} إلى الحمل.

% --- 3.2 إدخال بارامترات العناصر ---
\subsection{إدخال بارامترات العناصر \el{(Parameter Configuration)}}

تم ضبط قيم العناصر بناءً على الحسابات النظرية التي أجريناها سابقاً لتردد رنيني قدره \el{$150 Hz$}. توضح الأشكال التالية كيفية إدخال هذه القيم داخل بلوكات \el{Simulink}:

\begin{enumerate}
    \item \textbf{قيم الحثيات \el{(Inductances)}:} تم ضبط الملفات التسلسلية على القيمة المفروضة \el{$0.002 H$}.
    \begin{figure}[H]
        \centering
        \includegraphics[width=0.7\linewidth]{15 - Inductances of Low-Pass filter.png}
        \caption{نافذة بارامترات ملفات الخنق في النموذج المطور}
    \end{figure}

    \item \textbf{قيم المقاومات \el{(Resistances)}:} تم إدخال قيمة مقاومة التخميد المحسوبة \el{$0.628 \Omega$}.
    \begin{figure}[H]
        \centering
        \includegraphics[width=0.7\linewidth]{16 - Resistances of Low-Pass filter.png}
        \caption{إدخال قيم المقاومات الخاصة بفرع الترشيح}
    \end{figure}

    \item \textbf{قيم السعات \el{(Capacitances)}:} تم ضبط المكثفات على القيمة \el{$562.9 \mu F$}.
    \begin{figure}[H]
        \centering
        \includegraphics[width=0.7\linewidth]{17 - Capacitances of Low-Pass filter.png}
        \caption{ضبط قيم السعات لمرشح التمرير المنخفض}
    \end{figure}
\end{enumerate}

% --- القسم الرابع: مناقشة وتحليل النتائج ---
\section{مناقشة وتحليل النتائج \el{(Results \& Discussion)}}

بعد إتمام عملية النمذجة، نلاحظ الشكل النهائي للدارة في بيئة \el{Simulink} كما هو موضح أدناه:

\begin{figure}[H]
    \centering
    \includegraphics[width=1.05\linewidth]{20 - Circuit with Low-Pass Filter.png}
    \caption{المخطط الصندوقي الكامل للدارة بعد إضافة مرشح التمرير المنخفض}
\end{figure}

أظهرت نتائج المحاكاة تحسناً جذرياً في جودة القدرة، ويمكن تلخيص النتائج كما يلي:

\begin{itemize}
    \item \textbf{تحليل التيار:} نلاحظ انخفاض نسبة التشوه التوافقي الكلي للتيار \el{($THD_I$)} بشكل كبير جداً، حيث تخلص المرشح من كافة الرتب التوافقية العالية.
    \begin{figure}[H]
        \centering
        \includegraphics[width=0.9\linewidth]{18 - CurrentData with Low-Pass Filter.png}
        \caption{طيف التردد للتيار بوجود مرشح التمرير المنخفض}
    \end{figure}

    \item \textbf{تحليل الجهد:} أصبحت موجة الجهد أكثر نقاءً، واستقرت نسبة \el{$THD_V$} ضمن المعايير المسموح بها عالمياً.
    \begin{figure}[H]
        \centering
        \includegraphics[width=0.9\linewidth]{19 - VoltageData with Low-Pass Filter.png}
        \caption{طيف التردد للجهد بوجود مرشح التمرير المنخفض}
    \end{figure}
\end{itemize}
